% ==================== PARTIE A — VARIABLES ====================
% Étape de conception : modéliser les grandeurs de base d'une session de charge
% (puissance, durée, énergie, tarif, montant) et choisir les bons types.
%
% Progression :
%   A1  trace de valeurs (tableau à compléter)
%   A2  QCM — division entière et modulo (conversion min -> h)
%   A3  QCM — priorité des opérateurs (division exacte) + expression générale
%   A4a compléter les déclarations (types imposés)
%   A4b compléter les 3 expressions de calcul
%   (l'ancien QCM flottant est déplacé en partie F)

\element{var-types}{
    \noindent\libelleProchaineQuestion\\
    Pour chacune des grandeurs ci-dessous manipulées par la borne, cocher \textbf{une seule case} :
    le type le plus adapté pour la variable qui la représente.
    (\texttt{char} et \texttt{string} ne sont proposés ici qu'à titre d'options de reconnaissance.)

    \begin{center}
    {\sansNumeroQuestion\renewcommand{\arraystretch}{1.5}
    \begin{tabular}{|p{5.6cm}|p{10.6cm}|}
        \hline
        \textbf{Grandeur} & \textbf{Type C le plus adapté} \\ \hline
        Nombre de badges utilisateurs enregistrés dans la borne &
        \begin{question}{var-types-nbsessions}\begin{reponseshoriz}[o] \bonne{\texttt{int}} \mauvaise{\texttt{float}} \mauvaise{\texttt{char}} \mauvaise{\texttt{string}} \end{reponseshoriz}\end{question} \\ \hline
        Tension réseau mesurée (ex. 398,6 V) &
        \begin{question}{var-types-tension}\begin{reponseshoriz}[o] \mauvaise{\texttt{int}} \bonne{\texttt{float}} \mauvaise{\texttt{char}} \mauvaise{\texttt{string}} \end{reponseshoriz}\end{question} \\ \hline
        Nombre de coupures réseau détectées depuis la mise en service &
        \begin{question}{var-types-point}\begin{reponseshoriz}[o] \bonne{\texttt{int}} \mauvaise{\texttt{float}} \mauvaise{\texttt{char}} \mauvaise{\texttt{string}} \end{reponseshoriz}\end{question} \\ \hline
        Rendement du convertisseur de puissance (ex. 0,94) &
        \begin{question}{var-types-tarif}\begin{reponseshoriz}[o] \mauvaise{\texttt{int}} \bonne{\texttt{float}} \mauvaise{\texttt{char}} \mauvaise{\texttt{string}} \end{reponseshoriz}\end{question} \\ \hline
        Un caractère représentant la catégorie d'un utilisateur (ex. 'A', 'B', 'C') &
        \begin{question}{var-types-duree}\begin{reponseshoriz}[o] \mauvaise{\texttt{int}} \mauvaise{\texttt{float}} \bonne{\texttt{char}} \mauvaise{\texttt{string}} \end{reponseshoriz}\end{question} \\ \hline
        Index du compteur d'énergie depuis l'installation (ex. 18\,452,7 kWh) &
        \begin{question}{var-types-energie}\begin{reponseshoriz}[o] \mauvaise{\texttt{int}} \bonne{\texttt{float}} \mauvaise{\texttt{char}} \mauvaise{\texttt{string}} \end{reponseshoriz}\end{question} \\ \hline
    \end{tabular}}
    \end{center}
    \grilleCommeUneQuestion
}

\element{var-qcm-declaration}{
    \begin{question}{var-qcm-declaration}
        Quelle ligne déclare correctement une variable entière \lstinline[language=C]|nb_sessions|
        initialisée à zéro ?
        \begin{multicols}{3}
            \begin{reponses}
                \bonne{\lstinline[language=C]|int nb_sessions = 0;|}
                \mauvaise{\lstinline[language=C]|nb_sessions int = 0;|}
                \mauvaise{\lstinline[language=C]|int nb_sessions == 0;|}
                \mauvaise{\lstinline[language=C]|int 0 = nb_sessions;|}
                \mauvaise{\lstinline[language=C]|(nb_sessions == 0);|}
                \mauvaise{\lstinline[language=C]|int nb_sessions == 0;|}
            \end{reponses}
        \end{multicols}
    \end{question}
}

\element{var-trace}{
    \begin{question}{var-trace-1}
        Avant d'écrire le logiciel, on vérifie la façon dont la borne manipule ses grandeurs.
        Le fragment ci-dessous traite une session de 7 kW pendant 150 minutes,
        facturée 0,22 €/kWh.
        
        Compléter le tableau en donnant la valeur de la variable affectée par chacune des lignes
        repérées \texttt{L1} à \texttt{L5}, \textbf{en respectant le type de la variable}
        (une valeur entière s'écrit sans décimale).
        \evaluationProf[2]

        \begin{minipage}[t]{.56\textwidth}
            \lstinputlisting[language=c]{sources/codes/a1-trace-grandeurs.c}
        \end{minipage}\hfill
        \begin{minipage}[t]{.42\textwidth}~\\
            {\renewcommand{\arraystretch}{2}
            \begin{tabular}{|c|l|p{2.6cm}|}
                \hline
                \textbf{Ligne} & \textbf{Variable} & \textbf{Valeur} \\ \hline
                L1 & \lstinline[language=C]|duree_h| & \\ \hline
                L2 & \lstinline[language=C]|energie_kwh| & \\ \hline
                L3 & \lstinline[language=C]|energie_kwh_ent|  & \\ \hline
                L4 & \lstinline[language=C]|energie_kwh| & \\ \hline
                L5 & \lstinline[language=C]|cout_euros|  & \\ \hline
            \end{tabular}}

            \medskip
        \end{minipage}
    \end{question}
}

\element{var-qcm-nomVariables}{
    Le développeur a choisi des noms de variables qui ne sont pas très explicites.
    Il souhaite les renommer pour mieux refléter leur rôle dans le calcul du coût
    d'une session de charge.
    
    \begin{center}
        \begin{minipage}{.74\textwidth}
            \lstinputlisting[language=c]{sources/codes/a2-qcm-nomVariables.c}
        \end{minipage}
    \end{center}
    
    \begin{question}{var-qcm-nomVariables}
        Quel nom de variable est le plus approprié pour remplacer \lstinline[language=C]|a| ?
        \begin{multicols}{4}
            \begin{reponses}
                \bonne{\lstinline[language=C]|duree_min|}
                \mauvaise{\lstinline[language=C]|heures_entamees|}
                \mauvaise{\lstinline[language=C]|reste_min|}
                \mauvaise{\lstinline[language=C]|minutes_restant|}
            \end{reponses}
        \end{multicols}
    \end{question}
}

\element{var-qcm-conversion}{
    Pour le ticket de stationnement, la borne convertit la durée (min) en
    \emph{heures entamées} et en minutes restantes.

    \begin{center}
        \begin{minipage}{.74\textwidth}
            \lstinputlisting[language=c]{sources/codes/a2-qcm-conversion.c}
        \end{minipage}
    \end{center}

    \begin{question}{var-lecture-printf}
        Pour \lstinline[language=C]|heures_entamees|
        valant \texttt{2} et \lstinline[language=C]|reste_min| valant \texttt{7}, qu'affiche la dernière
        ligne (l'appel à \lstinline[language=C]|printf|) du programme ci-dessous ?
        \begin{multicols}{2}
            \begin{reponses}
                \bonne{\texttt{2 7}}
                \mauvaise{\texttt{7 2}}
                \mauvaise{\texttt{27}}
                \mauvaise{\texttt{722}}
                \mauvaise{\texttt{\%d \%d}}
                \mauvaise{\texttt{heures\_entamees reste\_min}}
            \end{reponses}
        \end{multicols}
    \end{question}
    


    \begin{question}{var-qcm-conversion}
        Qu'affiche ce fragment de code complet (avec \lstinline[language=C]|duree_min| égal à 185) ?
        \begin{multicols}{6}
            \begin{reponses}
                \bonne{\texttt{4 5}}
                \mauvaise{\texttt{3 5}}
                \mauvaise{\texttt{4 65}}
                \mauvaise{\texttt{3 65}}
                \mauvaise{\texttt{4 55}}
                \mauvaise{\texttt{3 55}}
            \end{reponses}
        \end{multicols}
    \end{question}
}

\element{var-qcm-priorite}{
    Le stationnement est majoré de 3~\euro{} par heure entamée \emph{au-delà} de
    120 minutes, en plus de 50 centimes de frais de service. 
    
    Un développeur inexpérimenté écrit :
    
    \begin{center}
        \begin{minipage}{.74\textwidth}
            \lstinputlisting[language=c]{sources/codes/a3-qcm-priorite.c}
        \end{minipage}
    \end{center}
    
    \begin{question}{var-qcm-priorite-180}
        Qu'affiche ce fragment de code si on entre une durée de 180 minutes ?
        \begin{multicols}{6}
            \begin{reponses}
                \bonne{350 cts}
                \mauvaise{50 cts}
                \mauvaise{200 cts}
                \mauvaise{950 cts}
                \mauvaise{0 cts}
                \mauvaise{300 cts}
            \end{reponses}
        \end{multicols}
    \end{question}

    \begin{question}{var-qcm-priorite-150}
        Qu'affiche ce fragment de code si on entre une durée de 150 minutes ?
        \begin{multicols}{6}
            \begin{reponses}
                \mauvaise{350 cts}
                \bonne{50 cts}
                \mauvaise{200 cts}
                \mauvaise{950 cts}
                \mauvaise{0 cts}
                \mauvaise{300 cts}
            \end{reponses}
        \end{multicols}
    \end{question}

    \begin{question}{var-qcm-priorite-corrige}
        
        Proposer une correction du programme pour qu'il calcule correctement le montant à facturer.

        \evaluationProfGrille{2}{2}
    \end{question}
}

\element{var-completion-decl}{
    On facture maintenant une session réelle. La durée est saisie \textbf{en minutes},
    le tarif est exprimé en €/kWh, et tout montant comprend 0,50 € de frais de service.
    Les types sont imposés : \lstinline[language=C]|puissance_kw| est un nombre à virgule
    (\lstinline[language=C]|float|), \lstinline[language=C]|duree_min| un nombre entier
    (\lstinline[language=C]|int|), \lstinline[language=C]|tarif_kWh| un nombre à virgule
    (\lstinline[language=C]|float|).

    \textbf{Rappel :} énergie (kWh) $=$ puissance (kW) $\times$ durée (h).

    \lstinputlisting[language=c,escapeinside={(*}{*)}]{sources/codes/a4-cout-session-a-completer.c}

    \begin{question}{var-completion-decl}

        \evaluationProf[1]

        Compléter \textbf{directement dans les trois premiers cadres pointillés} la déclaration
        des variables \lstinline[language=C]|puissance_kw|, \lstinline[language=C]|duree_min|
        et \lstinline[language=C]|tarif_kWh|.

    \end{question}
    \begin{question}{var-completion-calcul}
        \evaluationProf[3]

        Dans le programme de la question précédente, compléter les \textbf{trois expressions
        de calcul} manquantes (cadres pointillés de \lstinline[language=C]|duree_h|,
        \lstinline[language=C]|energie_kwh| et \lstinline[language=C]|montant_euros|).
        Attention au type du résultat de chaque calcul intermédiaire.

        Utiliser, si besoin, l'espace ci-dessous pour justifier un choix de type ou de conversion.

    \end{question}
}
